\documentclass[a4paper, 9pt]{article}
\usepackage{geometry}
%% % Geometry: es para modificar los márgenes del documento
\usepackage[left=3cm, right=2.5cm, top=2.5cm, bottom=2.5cm]{geometry}
%% lipsum: para generar texto de ejemplo. Se puede eliminar una vez que se eliminen todos los \lipsum[] del documento.
\usepackage{lipsum}
% lscape: en caso de querer rotar una hoja, buscar información en internet en caso de ser requerido. 
\usepackage{lscape}
\usepackage{float}
% inputenc: para que latex acepte caracteres latinos como los acentos y la letra ñ.
\usepackage[utf8]{inputenc}
% babel: para traducir los títulos que vienen originalmente en inglés. Ejemplo: Fecha, Chapter, Bibliography, Appendix, etc.
\usepackage[spanish,es-tabla]{babel}
% natbib: para poder citar utilizando paréntesis redondos con \citep{•} o sin parentesis con \cite{•}
\usepackage[numbers]{natbib}
\usepackage{float} % para usar [H] y obligar que las figuras o tablas aparezcan donde es requerido.
\usepackage[pdftex]{graphicx} % graphicx: para incorporar imágenes. Recordar que las imágenes 'gif' no son aceptadas por Latex, se sugiere utilizar formato png por su calidad, en segunda intantcia jpg.
\usepackage{parskip} % parskip: par no dejar sangrías e insertar espacios entre párrafos en su lugar.
\usepackage{amsmath} %paquete para escribir fórmulas matemáticas.
\usepackage{amsfonts}%paquete para escribir fórmulas matemáticas.
\usepackage{amssymb} %paquete para escribir fórmulas matemáticas.
\usepackage{amsbsy}
\usepackage{upgreek}
\usepackage[usenames,dvipsnames,svgnames,table]{xcolor} %xcolor: para definir colores y dar color a tablas.
\usepackage{multirow}

\graphicspath{ {figures/} }
\usepackage{array}

% hyperref: define opciones especiales para el documento PDF producido.
\usepackage[pdftex, bookmarksnumbered,  pagebackref, colorlinks=true, citecolor=DarkBlue, linkcolor=DarkBlue!30!Black, urlcolor=Black,bookmarksopen]{hyperref}

% El paquete fancyhdr es para definir opciones de encabezado y pié de página
% Según el formato existente a la fecha (agosto de 2015) esto no se considera
% Su utilización en este caso es para situar el número de página en la parte
% inferior derecha de la página.

\usepackage{fancyhdr} % activamos el paquete
	\pagestyle{fancy} % seleccionamos un estilo
	\lhead{} % texto izquierda de la cabecera
	\chead{} % texto centro de la cabecera
	\rhead{\textcolor[gray]{0.5}{\textit{\nouppercase \leftmark}}} % Nombre del capítulo. \nouppercase: uso de minúsculas
	\lfoot{} % texto izquierda del pie
	\cfoot{} % imagen centro del pie
	\rfoot{\textcolor[gray]{0.5}{\thepage}} % Número de página a la derecha, abajo
	\renewcommand{\headrulewidth}{0.2pt} % grosor de la línea de la cabecera

\fancypagestyle{detailed}{
    \fancyhf{} % clear all header and footers
    \fancyfoot[R]{\textcolor[gray]{0.5}{\thepage}}
	%\fancyhead{}    
    \renewcommand{\headrulewidth}{0pt}
 }
 
\usepackage{etoolbox}
\patchcmd{\chapter}{\thispagestyle{plain}}{\thispagestyle{detailed}}{}{}

%times: para uar letra tipo Times New Roman
\usepackage{times}

%separación entre líneas (1.2 espacios). En word es interlineado exacto a 12 pts
\renewcommand{\baselinestretch}{1.2} 

\usepackage{titlesec} % para poder modificar los títulos

% Para la numeración de tablas y figuras.
\renewcommand\thefigure{\arabic{section}.\arabic{figure}} % Genera numeración X.Y
\renewcommand\thetable{\arabic{section}.\arabic{table}} % Genera numeración X.Y
\numberwithin{figure}{section} %Hace que la primera figura de cada sección X sea X.1
\numberwithin{table}{section} %Hace que la primera tabla de cada sección X sea X

\usepackage{booktabs} % Para trabajar con opciones especiales de tablas.
\usepackage{caption}

\usepackage{setspace}

% El índice de tablas e imágenes se superpone el texto al número de la figura o tabla.
% Esta configuración arregla dicho problema, modificar el 3.0 de ser necesario.
\usepackage{tocloft}
\addtolength{\cftfignumwidth}{3.0em}
\renewcommand{\cftfigpresnum}{\figurename\ }
\addtolength{\cfttabnumwidth}{3.0em}
\renewcommand{\cfttabpresnum}{\tablename\ }
\usepackage[utf8]{inputenc}
\usepackage{amssymb}
\usepackage{graphicx}
\usepackage{lscape}
\usepackage{float}
\usepackage{changepage}
\usepackage{capt-of}
\usepackage{wrapfig}

\begin{document}

\begin{titlepage}
		\begin{center}
		\vspace{5 mm}\textbf{\Large UNIVERSIDAD POLITÉCNICA DE MADRID} \par
		\vspace{10 mm}
      			\begin{figure}[htb]
				\begin{center}
					\includegraphics[scale=0.3]{./images/logo-upm.eps}
				\end{center}
			\end{figure}\par
  
      %\vspace{10 mm}\textbf{\large \degree} \par
      %\vspace{10 mm}\textbf{\large \faculty} \par
      %\vspace{20 mm}\textbf{\large \department } \par


      \vspace{20 mm}\textsc{\Large Homework 2 \\ Modelo de regresión para los precios de diamantes } \par
      
      \vspace{10 mm}\textsc{STATISTICAL DATA ANALYSIS}\par
      \vspace{10 mm}\textsc{AUTORES: \\ Cristian Abrante Dorta, \\ Álvaro Arranz Domínguez, \\ Ángel González López, \\ Daniel Saiz González}\par
		
	\rule{80mm}{0.1mm}\\      
      
	  \vspace{10 mm}
      \textsc{Master's Programme in ICT Innovation: Data Science}\par 
      %\vfill\textrm{\address. \hfill \monthyear}
      
      
      \vspace{10 mm}\textsc{ESCUELA TÉCNICA SUPERIOR DE INGENIERÍA DE INFORMÁTICA}		
		\end{center}
\end{titlepage}


\newgeometry{textwidth=18cm,textheight=27cm}

\section{Descripción del conjunto de datos}
\label{sec:descripcion-datos}
\noindent

El conjunto de datos que estudiaremos se trata de una lista de 308 piedras preciosas. Este conjunto de datos se ha creado con fines educativos e incluye las siguientes variables:

\begin{itemize}
    \item \texttt{caratage}: Esta variable numérica continua se refiere a los quilates que posee el diamante, unidad de medida utilizada para pesar gemas y otras piedras preciosas. La correspondencia es de 1 quilate con 0.2 gramos.
    \item \texttt{purity}: Nivel de pureza del color del diamante. Varía desde el nivel \textbf{D} hasta el \textbf{I}, donde el \textbf{D} se corresponde con el nivel más puro.
    \item \texttt{clarity}: Nivel de claridad del diamante. Se corresponde con los valores: \textbf{IF, VVS1, VVS2, VS1, VS2}, ordenados de mejor a peor.
    \item \texttt{certificate}: Certificado de calidad emitido por una institución de renombre. Puede tener los siguientes valores: \textbf{GIA, IGI y HRD}.
    \item \texttt{price}: Precio que tiene el diamante, expresado en dólares de Singapur.
\end{itemize}

\section{Preguntas de investigación}
\noindent
El objetivo de esta investigación es dar respuesta al conjunto de preguntas formuladas aplicando diferentes técnicas de visualización y transformación de datos.

%%%%%%%%%%%%%%%%%%%%%%%%%%%%%%%%%%%%%%%%%
%   Pregunta 1
%%%%%%%%%%%%%%%%%%%%%%%%%%%%%%%%%%%%%%%%%
\subsection{Representa la gráfica \texttt{price} vs. \texttt{caratage} y la gráfica \texttt{log(price)} vs \texttt{caratage}. Decide que variable de respuesta es mejor utilizar.}
\label{subsec:question-1}



%%%%%%%%%%%%%%%%%%%%%%%%%%%%%%%%%%%%%%%%%
%   Pregunta 2
%%%%%%%%%%%%%%%%%%%%%%%%%%%%%%%%%%%%%%%%%
\subsection{Encuentra una manera adecuada de incluir, además de \texttt{caratage}, la otra información categórica disponible: \texttt{clarity}, \texttt{color} y \texttt{certificate}. Comenta el modelo que se ha ajustado, realizando un análisis básico de los residuos.}
\label{subsec:question-2}


%%%%%%%%%%%%%%%%%%%%%%%%%%%%%%%%%%%%%%%%%
%   Pregunta 3
%%%%%%%%%%%%%%%%%%%%%%%%%%%%%%%%%%%%%%%%%
\subsection{Crea dos acciones para remediar el resultado anterior}
\label{subsec:question-3}

\subsubsection{Crea una nueva variable categórica para segregar las piedras de acuerdo al valor de \texttt{caratage}.}

\subsubsection{Añade el cuadrado de \texttt{caratage} como una nueva variable explicatoria.}



%%%%%%%%%%%%%%%%%%%%%%%%%%%%%%%%%%%%%%%%%
%   Pregunta 4
%%%%%%%%%%%%%%%%%%%%%%%%%%%%%%%%%%%%%%%%%
\subsection{¿Cuál de las acciones anteriores prefieres y por qué? Justifica la respuesta en términos de interpretabilidad y validez de los supuestos.}
\label{subsec:question-4}



%%%%%%%%%%%%%%%%%%%%%%%%%%%%%%%%%%%%%%%%%
\begin{thebibliography}{9}
\bibitem{bortner1969short}
    Journal of chronic diseases (Bortner, Rayman W)
    \newblock {\em A short rating scale as a potential measure of pattern A behavior}, 1969.

\end{thebibliography}

\end{document}



